\documentclass[12pt]{article}
\usepackage[utf8]{inputenc}
\usepackage[T1]{fontenc}
\usepackage{hyperref}
\title{Kuku-Ugbanh Language Resources for K-12 Curriculum Development}
\author{Compiled by Claude Sonnet 5 (claude-sonnet-5) for Andr\'as Kornai}
\date{}
\begin{document}
\maketitle
\section{Introduction}
Very little was found. In the search results, Kuku-Ugbanh (ugb) appears as one of the named varieties or patrilects within the Kugu Nganhcara complex of Cape York, Australia. No wordlist, grammar or teaching material specific to Ugbanh was found.
\section{Primary Educational Resources}
None found.
\section{Linguistic and Reference Materials}
\begin{itemize}
\item Wikipedia article on the Kugu Nganhcara language, listing Ugbanh among its varieties and the ugb code: \url{https://en.wikipedia.org/wiki/Kugu_Nganhcara_language}
\item Endangered Languages Project entry, probably for the Kugu complex (the snippet was not checked): \url{https://endangeredlanguages.com/node/92026}
\end{itemize}
\section{Religious and Community Resources}
None found.
\section{Digital and Audio Resources}
None found for Ugbanh. A search suggested checking the AIATSIS and Pama Language Centre archives, but this was not done.
\section{Curriculum Development Context}
Any revitalization work would probably have to draw on materials for the wider Kugu Nganhcara complex and on archival sources. Those were not located in this search.
\end{document}